\chapter{Declaración de uso de inteligencia artificial generativa}
\label{cap:anexoh}

\section*{H.1. Declaración}

Durante el desarrollo del TFM se utilizaron asistentes de inteligencia artificial
generativa como herramientas auxiliares de programación, consulta y revisión textual.
Su uso se realizó bajo supervisión de los autores y no sustituyó las decisiones
metodológicas, la ejecución de experimentos, la comprobación de resultados ni la
responsabilidad sobre la versión final de la memoria.

\section*{H.2. Herramientas, finalidad y control}

\begin{table}[htbp]
\centering
\caption{Uso declarado de herramientas de IA generativa}
\label{tab:anexoh-herramientas}
\small
\begin{tabularx}{\textwidth}{|p{3.1cm}|X|p{4.0cm}|}
\hline
\textbf{Herramienta} & \textbf{Uso auxiliar} & \textbf{Control aplicado} \\ \hline
GitHub Copilot \cite{GitHubCopilot2023} & Completado de código, apoyo en depuración y revisión de comentarios técnicos. & Revisión manual, ejecución de pruebas y aceptación expresa de cada cambio. \\ \hline
Asistentes basados en LLM, incluido Google Gemini \cite{GoogleGemini2026} & Consulta técnica, contraste de alternativas, corrección lingüística y apoyo en LaTeX. & Verificación contra código, resultados experimentales y fuentes primarias. \\ \hline
Herramientas de búsqueda asistida & Localización inicial de bibliografía y normativa oficial. & Validación posterior de autoría, fecha, DOI y procedencia oficial. \\ \hline
\end{tabularx}
\end{table}

El equipo utilizó además un enfoque de desarrollo guiado por especificaciones para
coordinar tareas y conservar trazabilidad. Esta organización no genera resultados
académicos por sí misma: las métricas incluidas en la memoria proceden de ejecuciones
reproducibles y fueron revisadas por los autores.

\section*{H.3. Límites y responsabilidad}

Las herramientas generativas no asignaron las etiquetas P1--P4 finales, no eligieron el
modelo de Capa 2 y no determinaron las conclusiones del estudio. La auditoría del
conjunto de datos, el diseño de supervisión débil, el entrenamiento, la evaluación y el
análisis de errores fueron ejecutados mediante procedimientos controlados por el equipo.

Los autores asumen la responsabilidad sobre el código, la selección de fuentes, los
resultados numéricos y la redacción final. El historial del repositorio conserva la
trazabilidad del desarrollo:

\begin{center}
\url{https://github.com/jalbfil/tfe_muia_unir_188}
\end{center}
